% Сборка из каталога проекта:
% latexmk -xelatex -interaction=nonstopmode -halt-on-error analog_power_lm317_lm337.tex
\documentclass[a4paper,landscape,10pt]{article}
\PassOptionsToPackage{dvipsnames}{xcolor}
\usepackage{iftex}

\ifTUTeX
\usepackage{xltxtra}

\usepackage{fontspec}
\usepackage{polyglossia}
\setmainlanguage[babelshorthands=true]{russian}
\setotherlanguage{english}

\setmainfont{Times New Roman}
\setromanfont{Times New Roman} 
\setsansfont{Arial} 
\setmonofont{Courier New} 

\newfontfamily{\cyrillicfont}{Times New Roman} 
\newfontfamily{\cyrillicfontrm}{Times New Roman}
\newfontfamily{\cyrillicfonttt}{Courier New}
\newfontfamily{\cyrillicfontsf}{Arial}

% \usepackage{unicode-math}
\usepackage{amssymb}
\usepackage{amsfonts}

\else
% For pdfTeX we must use old
% 8-bit font technologies
\usepackage[T2A]{fontenc}
\usepackage[english,russian]{babel}
\usepackage[utf8]{inputenc}
\usepackage[english, russian]{babel}

%\usepackage{newtxmath}
\usepackage{amsmath}
\usepackage{amsthm}
\usepackage{amssymb}
\usepackage{amsfonts}
%\usepackage{mathtext}
\fi

%Hyphenation rules
\usepackage{hyphenat}
\hyphenation{ма-те-ма-ти-ка вос-ста-нав-ли-вать ото-бра-жа-ют-ся сге-не-ри-ро-ван-ном сим-во-лы}
% \usepackage[english, russian]{babel}

\usepackage{graphicx}
\graphicspath{ {docs/images/} {../images/} {./images/} }

\usepackage{tikz}
\usepackage{pgfplots} % для построения графиков
\usepackage{tkz-euclide}
\pgfplotsset{compat=1.18} % версия для совместимости
\usepackage{empheq}
\usetikzlibrary{positioning, calc, intersections}

\usepackage{float}

% \usepackage{comment}

\usepackage[hidelinks]{hyperref}

\newtheorem{theorem}{Теорема}[section]
\newtheorem{corollary}{Corollary}[theorem]
\newtheorem{lemma}[theorem]{Lemma}

\addto\captionsrussian{%
  \renewcommand{\figurename}{Рис.}%
  \renewcommand{\tablename}{Табл.}%
}

\usepackage{verbatim}

\newcommand{\cg}{\textnormal{\textsl{g}}}

\usepackage[dvipsnames]{xcolor}

\usepackage[siunitx, RPvoltages, american]{circuitikzgit}
\ctikzset{
    amplifiers/fill=cyan,
    sources/fill=green!20,
    diodes/fill=white,
    resistors/fill=violet,
    grounds/scale=1.2,
}

\usepackage[margin=12mm]{geometry}
\pagestyle{empty}
\setlength{\parindent}{0pt}
\ctikzset{resistors/fill=white,diodes/fill=white,grounds/scale=0.65,
  bipoles/length=1.05cm}
\newcommand{\uF}{\,\text{мкФ}}
\newcommand{\nF}{\,\text{нФ}}
\newcommand{\Ohm}{\,\text{Ом}}
\newcommand{\kOhm}{\,\text{кОм}}
\newcommand{\gnd}{\textnormal{GND}}

% Знак у плоского символа C показывает положительную обкладку электролита.
% #1: 1 = LM317D2T, 0 = LM337D2T. D2PAK, по даташитам onsemi.
% Топология одинакова; полярности и номера силовых выводов различаются.
\newcommand{\regulator}[1]{%
\begingroup
\ifnum#1=1
  \def\ic{LM317D2T}\def\uid{U1}\def\innum{3}\def\outnum{2}
  \def\inputnet{+V_{\!F}}\def\outputnet{+12}
  \def\ci{5}\def\cb{6}\def\co{9}\def\cc{11}\def\ca{13}
  \def\ra{3}\def\rb{4}\def\rv{1}\def\da{5}\def\db{6}\def\dc{9}
\else
  \def\ic{LM337D2T}\def\uid{U2}\def\innum{2}\def\outnum{3}
  \def\inputnet{-V_{\!F}}\def\outputnet{-12}
  \def\ci{7}\def\cb{8}\def\co{10}\def\cc{12}\def\ca{14}
  \def\ra{5}\def\rb{6}\def\rv{2}\def\da{7}\def\db{8}\def\dc{10}
\fi
\begin{circuitikz}[y=0.8cm,font=\small,line width=0.6pt]
  \path[use as bounding box] (-0.8,-4.9) rectangle (25,2.3);
  \draw[fill=cyan!6] (6.5,-0.7) rectangle (9.5,0.7);
  \node at (8,0.15) {\uid\quad\ic};
  \node[font=\scriptsize,anchor=west] at (6.65,-0.3) {IN};
  \node[font=\scriptsize,anchor=east] at (9.35,-0.3) {OUT};
  \node[font=\scriptsize,anchor=south] at (8,-0.7) {ADJ};
  \draw (0,0) node[ocirc] {} node[above right] {$\inputnet$}
    -- (6.5,0) node[above left,font=\footnotesize] {\innum};
  \draw (9.5,0) node[above right,font=\footnotesize] {\outnum}
    -- (24,0) node[ocirc] {} node[above] {$\outputnet$ В};
  \draw (8,-0.7) node[below left,font=\footnotesize] {1}
    -- (8,-1.65) -- (15.7,-1.65);
  \draw (11,0) to[R,*-*] (11,-1.65)
    to[R,l_=$R_{\rb}\;820\Ohm$] (11,-3.1)
    to[vR,l_=$RV_{\rv}\;200\Ohm$] (11,-4.6) node[ground] {};
  \node[anchor=east] at (10.75,-1.1) {$R_{\ra}\;110\Ohm$};
  \draw (1.5,0) to[C,l_=$C_{\cb}\;100\nF$,*-]
    (1.5,-2.6) node[ground] {};
  \draw (4.7,0) to[C,l_=$C_{\ci}\;10\uF/50\text{ В}$,*-]
    (4.7,-2.6) node[ground] {};
  \draw (15.7,-1.65) to[C,l_=$C_{\ca}\;10\uF/25\text{ В}$]
    (15.7,-4.3) node[ground] {};
  \draw (18.2,0) to[C,l_=$C_{\co}\;10\uF/25\text{ В}$,*-]
    (18.2,-2.6) node[ground] {};
  \draw (21,0) to[C,l_=$C_{\cc}\;100\nF$,*-]
    (21,-2.6) node[ground] {};
  \draw (24,-2.6) node[ocirc] {} -- (23.6,-2.6) node[ground] {};
  \node[above] at (24,-2.6) {GND};
  \ifnum#1=1
    % VD5: OUT -> IN; VD6: ADJ -> OUT.
    \draw (11,0) -- (11,1.6) to[D,l_=$VD_{\da}$] (5.5,1.6)
      to[short,-*] (5.5,0);
    \draw (13.4,-1.65) to[D,l_=$VD_{\db}$,*-*] (13.4,0);
    % VD9 ограничивает отрицательный выброс на положительной шине.
    \draw (22.6,-2.6) node[ground] {}
      to[D,l_=$VD_{\dc}$] (22.6,0) node[circ] {};
    \node at (4.3,-0.8) {$+$};
    \node at (15.3,-2.25) {$+$};
    \node at (17.8,-0.8) {$+$};
  \else
    % VD7: IN -> OUT; VD8: OUT -> ADJ.
    \draw (5.5,0) to[short,*-] (5.5,1.6)
      to[D,l=$VD_{\da}$] (11,1.6) -- (11,0);
    \draw (13.4,0) to[D,l=$VD_{\db}$,*-*] (13.4,-1.65);
    % VD10 ограничивает положительный выброс на отрицательной шине.
    \draw (22.6,0) to[D,l=$VD_{\dc}$,*-]
      (22.6,-2.6) node[ground] {};
    \node at (4.3,-1.85) {$+$};
    \node at (15.3,-3.7) {$+$};
    \node at (17.8,-1.85) {$+$};
  \fi
\end{circuitikz}%
\endgroup}

\begin{document}
{\LARGE\bfseries Аналоговое питание ±12 В для фонокорректора}\par
\vspace{2mm}
\textbf{Лист 1. Вторичные обмотки, двухполупериодный выпрямитель и CRC-фильтры.}
Нагрузка: до 60 мА на плечо, сверх тока задающих делителей стабилизаторов.

\begin{center}
\resizebox{0.96\linewidth}{!}{%
\begin{circuitikz}[y=0.85cm,font=\small,line width=0.6pt]
  \path[use as bounding box] (-2.8,-5.3) rectangle (23.5,5.3);
  % Две вторичные секции; первичная цепь не входит в эту схему.
  \node[align=center,font=\bfseries] at (-0.3,4.8)
    {T1 — штатный трансформатор\\«Мелодия-103М»};
  \draw (-1.1,3.4) to[L] (-1.1,1.1);
  \draw (-1.1,-1.1) to[L] (-1.1,-3.4);
  \draw (-1.75,-3.5) -- (-1.75,3.5);
  \draw (-1.9,-3.5) -- (-1.9,3.5);
  \draw (-1.1,3.4) -- (0.3,3.4) node[above] {5}
    to[fuse,l=$F1$] (2.7,3.4) -- (7,3.4) -- (7,2);
  \draw (-1.1,1.1) -- (0.3,1.1) node[above] {6} -- (1.4,1.1);
  \draw (-1.1,-1.1) -- (0.3,-1.1) node[below] {$6'$} -- (1.4,-1.1);
  \draw (1.4,1.1) -- (1.4,-1.1);
  \draw (1.4,0) node[circ] {} -- (2.4,0) node[ground] {};
  \node[above right] at (1.4,0) {CT / GND};
  \node[font=\footnotesize,anchor=west] at (-0.7,2.15) {$14\text{–}15$ В$\sim$};
  \node[font=\footnotesize,anchor=west] at (-0.7,-2.15) {$14\text{–}15$ В$\sim$};
  \fill (-1.35,3.1) circle (1.5pt);
  \fill (-1.35,-1.4) circle (1.5pt);
  % Нижний AC-провод пересекает отрицательный DC-провод без соединения.
  \node[jump crossing,/tikz/circuitikz/bipoles/crossing/size=0.4]
    (X) at (3.9,-3.4) {};
  \draw (-1.1,-3.4) -- (0.3,-3.4) node[below] {$5'$}
    to[fuse,l_=$F2$] (2.7,-3.4) -- (X.west);
  \draw (X.east) -- (7,-3.4) -- (7,-2);
  % Один мост из четырёх диодов. Верх/низ = AC, право = +, лево = -.
  \draw (7,2) to[D,l=$VD1$] (9,0);
  \draw (7,-2) to[D,l_=$VD2$] (9,0);
  \draw (5,0) to[D,l=$VD3$] (7,2);
  \draw (5,0) to[D,l_=$VD4$] (7,-2);
  \node[circ] at (7,2) {};
  \node[circ] at (7,-2) {};
  \node[above right] at (7,2) {$\sim$};
  \node[below right] at (7,-2) {$\sim$};
  \node[above right] at (9,0) {$+$};
  \node[above left] at (5,0) {$-$};
  \draw (9,0) node[circ] {} -- (10.3,0) -- (10.3,4.0)
    -- (14,4.0) to[R,l=$R1\;22\Ohm/1\text{ Вт}$,*-*]
    (18.1,4.0) -- (23,4.0) node[ocirc] {} node[above] {$+V_{\!F}$};
  \draw (5,0) node[circ] {} -- (3.9,0) -- (X.north);
  \draw (X.south) -- (3.9,-4.0)
    -- (14,-4.0) to[R,l_=$R2\;22\Ohm/1\text{ Вт}$,*-*]
    (18.1,-4.0) -- (23,-4.0) node[ocirc] {} node[below] {$-V_{\!F}$};
  \draw (14,4.0) to[C,l_=$C1\;1000\uF/50\text{ В}$] (14,0);
  \draw (14,0) to[C,l_=$C2\;1000\uF/50\text{ В}$] (14,-4.0);
  \draw (18.1,4.0) to[C,l_=$C3\;470\uF/50\text{ В}$] (18.1,0);
  \draw (18.1,0) to[C,l_=$C4\;470\uF/50\text{ В}$] (18.1,-4.0);
  \draw (13,0) -- (23,0) node[ocirc] {} node[above] {GND};
  \node[circ] at (14,0) {};
  \node[circ] at (18.1,0) {};
  \node[ground] at (13,0) {};
  \node at (13.6,2.8) {$+$};
  \node at (13.6,-1.1) {$+$};
  \node at (17.7,2.8) {$+$};
  \node at (17.7,-1.1) {$+$};
  \node[font=\footnotesize,align=center] at (7,-4.9)
    {VD1–VD4: 1N4007\\F1, F2: T250 мА};
  \node[font=\footnotesize,align=center] at (13.8,4.9)
    {До фильтров: около $+19\text{–}21$ В / $-19\text{–}21$ В};
  \node[font=\footnotesize,align=center] at (21,2.1)
    {На вход U1\\(лист 2)};
  \node[font=\footnotesize,align=center] at (21,-2.1)
    {На вход U2\\(лист 2)};
\end{circuitikz}}
\end{center}

\begin{minipage}[t]{0.48\linewidth}
\textbf{Обмотки и фазировка}\par\small
Номера взяты из \texttt{docs/POWER.md}; указанные там напряжения — предварительные.
Для этого варианта измерить $14\text{–}15$ В между 5–6 и между $5'\text{–}6'$.
При соединении $6\text{–}6'$ между 5 и $5'$ должно быть $28\text{–}30$ В.
Точки у секций показывают \textbf{требуемую} фазировку; заводская маркировка не подтверждена.
Если напряжения вычитаются, соединение концов нужно изменить по результатам проверки.
Общий провод всегда берётся с фактической средней точки двух секций.
\par\vspace{1mm}
Секции сначала отсоединить от неисправного выпрямителя У5 и его перемычек.
Первичная обмотка, сетевые цепи, двигатель и прочие вторичные секции здесь не показаны.
\end{minipage}\hfill
\begin{minipage}[t]{0.48\linewidth}
\textbf{Принятые номиналы}\par\small
$C1$, $C2$: алюминиевые электролиты 1000 мкФ, 50 В, 105\,°C;
$C3$, $C4$: 470 мкФ, 50 В. Знаки «+» показывают положительные обкладки.
$R1$, $R2$: 22 Ом, 1 Вт, огнестойкие. F1 и F2 — плавкие предохранители
250 мА с задержкой; проверить запуск с выбранными конденсаторами.
\par\vspace{1mm}
При токе нагрузки 60 мА плюс около 11.4 мА делителя падение на каждом $22\Ohm$
составляет около 1.57 В. Частота пульсаций — 100 Гц при сети 50 Гц;
постоянная времени CRC-фильтра около 10.3 мс.
Если секции дают только $10\text{–}10.5$ В, для LM317/LM337 на ±12 В запаса недостаточно.
Выбор других секций требует сверки напряжений и нагрева.
\end{minipage}

\newpage
{\LARGE\bfseries Стабилизаторы LM317D2T/LM337D2T на ±12 В}\par
\vspace{2mm}
\textbf{Лист 2.} $+V_F$, GND и $-V_F$ приходят с листа 1.
Все одноимённые точки GND электрически соединены.
\begin{center}
\resizebox{0.90\linewidth}{!}{\regulator{1}}\par
\vspace{1mm}
\resizebox{0.90\linewidth}{!}{\regulator{0}}
\end{center}

\begin{minipage}[t]{0.48\linewidth}
\small
\textbf{Регулировка и обвязка}\par
$R3$, $R5$: 110 Ом, 1\%, 0.25 Вт; $R4$, $R6$: 820 Ом, 1\%, 0.25 Вт.
$RV1$, $RV2$: многооборотные подстроечники 200 Ом, включённые реостатами;
ползунок соединить с используемым крайним выводом.
До подключения фонокорректора выставить $+12.0$ и $-12.0$ В с контрольными
нагрузками по 1 кОм / 0.25 Вт от каждой выходной шины к GND; проверить и холостой ход.
Делители потребляют около $1.25/110=11.4$ мА и обеспечивают минимальную нагрузку 10 мА.
\par\vspace{1mm}
$C5$, $C7$: 10 мкФ / 50 В; $C9$, $C10$, $C13$, $C14$: 10 мкФ / 25 В —
алюминиевые электролиты. $C6$, $C8$, $C11$, $C12$: 100 нФ, от 50 В, керамика.
Входные и выходные конденсаторы устанавливать у стабилизаторов.
Выходной электролит $C10$ необходим для устойчивости LM337;
для него использовать обычный алюминиевый электролит, не сверхнизкоомный.
\par\vspace{1mm}
Распиновка D²PAK: U1 — ADJ 1, OUT 2, IN 3; U2 — ADJ 1, IN 2, OUT 3.
Теплоотводящая площадка каждого корпуса соединена с выводом 2:
у U1 это $+12$ В, у U2 — $-V_F$.
Полигоны теплоотвода раздельные, с GND не соединять; держать вдали от входа головки.
\end{minipage}\hfill
\begin{minipage}[t]{0.48\linewidth}
\small
\textbf{Защита, монтаж и проверка}\par
VD5–VD10: 1N4007. VD5/VD7 дают путь разряда выходных конденсаторов при падении входа;
VD6/VD8 — конденсаторов ADJ; VD9/VD10 ограничивают обратную полярность выходных шин.
\par\vspace{1mm}
Под нагрузкой минимальное $|V_F|$ желательно не ниже 15.5 В.
При 60 мА нагрузки суммарный ток каждого стабилизатора — около 71.4 мА.
Нагрев считать \textbf{вместе с делителем}: при $|V_F|=19$ В это около 0.50 Вт,
при 21 В — около 0.64 Вт на корпус.
По табличному $\theta_{JA}=70$ °C/Вт для D²PAK подъём температуры кристалла
оценивается в $35\text{–}45$ °C.
Результат зависит от платы: предусмотреть медь для теплоотвода на площадках корпусов
и проверить нагрев в корпусе при максимальной нагрузке и напряжении сети.
\par\vspace{1mm}
Возвраты выпрямителя вести к общей точке C1/C2;
возвраты нагрузки и делителей — отдельно от импульсов зарядки этих конденсаторов.
\end{minipage}
\vfill
\footnotesize
Выходы подключаются к шинам ±12 В фонокорректора; развязка непосредственно у NE5532 сохраняется.
Расчётный проект, без испытания макета.\par
Документация onsemi:
\href{https://www.onsemi.com/pdf/datasheet/lm317-d.pdf}{LM317/D},
\href{https://www.onsemi.com/pdf/datasheet/lm337-d.pdf}{LM337/D}
(локально: \texttt{docs/datasheet/ldo/LM337-D.PDF}).
\end{document}
